% !PRINTED 213
\weberpage{213}
%<<CONT
of $\varphi(\xi)$, just as the corresponding functions
$F_0(x), F_1(x) \cdots F_{n-2}(x)$ are formed from $x$ and the coefficients of $f(x)$.

If now, on the other hand, we carry out in $F(t, x)$, simultaneously with the
substitution (1), the substitution

\weq{4}{t = \frac{\alpha\tau+\beta}{\gamma\tau+\delta}}

then we obtain

\weq{5}{(\gamma\tau+\delta)^n f(t) = \varphi(\tau),}
\weq{6}{(\gamma\xi+\delta)(\gamma\tau+\delta)(t-x) = r(\tau-\xi)}

and, by forming from (5) the derivative, most simply by
differentiation by means of the formula

\[
\frac{dt}{d\tau} = \frac{r}{(\gamma\tau+\delta)^2}
\]

we obtain

\weq{7}{r(\gamma\tau+\delta)^{n-1}\frac{1}{n}f'(t) = (\gamma\tau+\delta)\frac{1}{n}\varphi'(\tau) - \gamma\varphi(\tau).}

From (5), (6) and (7) it follows that

\[
r(\gamma\tau+\delta)^{n-2}\frac{f(t)}{t-x} = \frac{\varphi(\tau)}{\tau-\xi}\frac{\gamma\xi+\delta}{\gamma\tau+\delta}
\]

\[
r(\gamma\tau+\delta)^{n-2}\frac{1}{n}f'(t) = -\frac{\gamma(\tau-\xi)\varphi(\tau)}{(\gamma\tau+\delta)(\tau-\xi)} + \frac{1}{n}\varphi'(\tau),
\]

and by subtraction of the two formulas

\[
r(\gamma\tau+\delta)^{n-2}\left(\frac{f(t)}{t-x} - \frac{1}{n}f'(t)\right) = \frac{\varphi(\tau)}{\tau-\xi} - \frac{1}{n}\varphi'(\tau),
\]

hence

\weq{8}{r(\gamma\tau+\delta)^{n-2}F(t, x) = \Phi(\tau, \xi)}

or, written more explicitly,

\weq{9}{\tau^{n-2}\Phi_0 + \tau^{n-3}\Phi_1 + \cdots + \Phi_{n-2} = r[(\alpha\tau+\beta)^{n-2}F_0 +(\alpha\tau+\beta)^{n-3}(\gamma\tau+\delta)F_1 +\cdots+(\gamma\tau+\delta)^{n-2}F_{n-2}].}

We now set, replacing in $F(t, x)$ the powers $t^k$ by
arbitrary variables $t_k$, and likewise in $\Phi(\tau, k)$ % [sic] printed \textquotedbl{}$\Phi(\tau,k)$\textquotedbl{}, not $\Phi(\tau,\xi)$
$\tau^k$ by $\tau_k$,

\weq{10}{\begin{aligned}
Y(t, x) &= t_{n-2}F_0 + t_{n-3}F_1 + \cdots + t_0 F_{n-2},\\
H(\tau, \xi) &= \tau_{n-2}\Phi_0 + \tau_{n-3}\Phi_1 + \cdots + \tau_0\Phi_{n-2},
\end{aligned}}

so that the substitutions of the Tschirnhausen transformation

\weq{11}{y = Y(t, x), \qquad \eta = H(\tau, \xi)}

read [\S.\,68, (12)].
%CONT>>
